\documentclass[12pt,a4paper]{article}
\usepackage[utf8]{inputenc}
\usepackage[T1]{fontenc}
\usepackage[brazil]{babel}
\usepackage[margin=2.4cm]{geometry}
\usepackage{parskip}
\usepackage{hyperref}
\usepackage{enumitem}
\usepackage{longtable}
\usepackage{array}
\usepackage{xcolor}

\hypersetup{
  colorlinks=true,
  linkcolor=blue,
  urlcolor=blue,
  citecolor=blue
}

\title{Relatório de Desenvolvimento do Protótipo Funcional\\\large 3\textsuperscript{o} Trabalho de Introdução ao Desenvolvimento de Jogos}
\author{João Pedro Araújo \and Túlio Mota Lima \and Rafael Luiz \and Otávio}
\date{31/07}

\begin{document}
\maketitle

\section*{Identificação}
\textbf{Jogo:} Music Guessr\\
\textbf{Disciplina:} Introdução ao Desenvolvimento de Jogos\\
\textbf{Entrega:} Protótipo funcional + relatório de desenvolvimento

\section*{1. Visão geral do protótipo}
O protótipo é um jogo musical web em que o jogador ouve trechos curtos de músicas e digita a resposta para pontuar. O sistema usa autenticação, ranking, painel administrativo e um fluxo de jogo baseado em rodada, tentativa e pontuação.

O projeto foi implementado como aplicação web, com frontend em React + Vite e backend em Node.js + Express + Prisma, usando PostgreSQL para persistência.

\section*{2. Objetivo de desenvolvimento}
Construir um protótipo funcional que entregasse:
\begin{itemize}[leftmargin=1.2em]
  \item interação lúdica significativa;
  \item core loop fluido;
  \item feedback visual e sonoro imediato;
  \item estabilidade técnica;
  \item ranking e persistência de desempenho;
  \item área administrativa para criar e organizar desafios.
\end{itemize}

\section*{3. Tecnologias utilizadas}
\subsection*{Frontend}
\begin{itemize}[leftmargin=1.2em]
  \item React + Vite
  \item React Router
  \item Motion / Framer Motion
  \item Tailwind CSS
  \item Axios
\end{itemize}

\subsection*{Backend}
\begin{itemize}[leftmargin=1.2em]
  \item Node.js + Express
  \item Prisma ORM
  \item PostgreSQL
  \item JWT para autenticação
  \item OAuth Google
  \item Integração com Deezer para busca e preview de músicas
\end{itemize}

\section*{4. Arquitetura do protótipo}
\subsection*{4.1 Frontend}
O frontend concentra as telas de login, registro, home, seleção de desafio, partida, ranking, perfil e área administrativa. A interface usa componentes reutilizáveis, animações de feedback e tema claro/escuro.

\subsection*{4.2 Backend}
O backend organiza autenticação, desafios, sessão de jogo, ranking, perfil e estatísticas. O jogo usa sessões em memória para controlar rodada atual, score e feedback durante a partida.

\subsection*{4.3 Banco de dados}
O modelo de dados inclui:
\begin{itemize}[leftmargin=1.2em]
  \item \texttt{Usuario}
  \item \texttt{Desafio}
  \item \texttt{DesafioMusica}
  \item \texttt{Tentativa}
\end{itemize}

Esses modelos sustentam usuários, desafios, músicas vinculadas e histórico de pontuação.

\section*{5. Como o jogo funciona}
\begin{enumerate}[leftmargin=1.4em]
  \item O jogador entra no sistema e faz login ou cadastro.
  \item Escolhe um desafio musical disponível.
  \item O sistema carrega o preview da música da rodada.
  \item O jogador digita o palpite.
  \item O backend valida a resposta com regras que toleram pequenas variações.
  \item Se acertar, recebe pontuação base + bônus por velocidade.
  \item O jogo avança para a próxima rodada.
  \item Ao final, as tentativas são persistidas e o ranking é atualizado.
\end{enumerate}

\section*{6. Critérios do professor e atendimento no protótipo}
\subsection*{6.1 Interação lúdica significativa e arquitetura de código}
Atendido. Cada ação do jogador gera retorno claro e imediato: envio de palpite, feedback de acerto ou erro, atualização da pontuação, avanço de rodada e finalização da sessão.

A arquitetura está separada em rotas, controllers, hooks e componentes reutilizáveis.

\subsection*{6.2 Calibração do flow e core loop}
Atendido. O ciclo principal está implementado e operacional:
\begin{quote}
Ouvir trecho $\rightarrow$ responder $\rightarrow$ receber feedback $\rightarrow$ avançar $\rightarrow$ finalizar sessão.
\end{quote}

A dificuldade altera a janela do trecho, o ritmo da partida e a pontuação, ajudando a manter o jogador engajado.

\subsection*{6.3 Integração estética e sonoplastia}
Atendido. O protótipo possui identidade visual consistente, animações de feedback, sinais sonoros de acerto, erro e conclusão, e interface que reforça a leitura do estado do jogo.

\subsection*{6.4 Storytelling e UX}
Atendido. A experiência é construída em torno da temática musical e da progressão por desafios. A interface permite leitura rápida do estado da partida e acesso direto às funções principais, sem excesso de fricção.

\subsection*{6.5 Estabilidade técnica, versionamento e QA}
Atendido em nível de protótipo web. O projeto possui build de frontend, backend funcional, autenticação, tratamento de erros em fluxos principais e repositório Git com histórico de desenvolvimento.

\section*{7. Decisões tomadas durante o desenvolvimento}
\begin{itemize}[leftmargin=1.2em]
  \item Separação entre frontend e backend para organização do código.
  \item Uso de Prisma + PostgreSQL para persistência relacional.
  \item Sessão de jogo em memória para simplificar o protótipo.
  \item Validação inteligente de respostas para reduzir frustração do jogador.
  \item Pontuação com bônus de tempo para reforçar o ritmo do jogo.
  \item Persistência apenas ao final da partida para manter o ranking coerente.
  \item Painel admin para criar e organizar desafios.
\end{itemize}

\section*{8. Participação dos membros}
\subsection*{João Pedro Araújo}
Backend, autenticação, sessão de jogo e ranking.

\subsection*{Rafael Luiz}
Modelagem de dados, Prisma, persistência e integração de estatísticas.

\subsection*{Túlio Mota Lima}
Frontend, interface, fluxo visual e feedback de interação.

\subsection*{Otávio}
Telas e navegação, ajustes de usabilidade e refinamento visual.

\section*{9. Entregáveis}
\begin{itemize}[leftmargin=1.2em]
  \item Relatório do desenvolvimento do protótipo.
  \item Link do GitHub.
  \item Código-fonte completo.
  \item Protótipo funcional em execução web.
\end{itemize}

\section*{10. Conclusão}
O protótipo entregue demonstra o core loop do jogo, a integração entre interface e backend, a persistência do desempenho e a possibilidade de evolução para novas funcionalidades. O projeto atende ao objetivo de apresentar um jogo funcional, jogável e coerente com os critérios do trabalho.

\section*{Observação para submissão}
Este relatório foi preparado para conversão em PDF. A versão em \texttt{.tex} está pronta para compilação.

\end{document}
